% 5.5 — Résultats
\section{Résultats}
\label{sec:resultats}

Le tableau~\ref{tab:resultats} rassemble les mesures de l'outil au commit de référence ;
les démonstrations qui permettent de les observer sont réunies en annexe~\ref{ann:captures}.

\begin{table}[htbp]
\centering
\caption{Résultats chiffrés au commit de référence \commitref{}}
\label{tab:resultats}
\small
\begin{tabularx}{\linewidth}{@{}L{6.4cm}Y@{}}
\toprule
\textbf{Mesure} & \textbf{Valeur} \\
\midrule
Code source / tests & $\approx$ \lignessrc{} / $\approx$ \lignestests{} lignes de Python \\
Fichiers de test / fonctions de test & \fichierstests{} / \fonctionstests{} (avant
  paramétrisation) \\
Table de règles & \nbregles{} règles, \nbexemples{} exemples exécutés comme tests,
  \nboutils{} outils, \nbfamilles{} familles à la demande \\
Règles dont la grammaire n'est pas vérifiée sur l'outil & \nbnonverifiees{} \\
Surfaces distantes avec canal ask/deny & \nbsurfaces{} \\
Suite complète, Linux & \testslinux \\
Suite complète, Windows & \testswindows \\
Portail de décision (\nbcorpus{} commandes) & Windows : 0 ligne différente ; Linux : une
  seule ligne, qui dépend des fichiers présents sur la machine \\
Corpus élargi (validation de M2.3b) & \nbcorpuslarge{} commandes rejouées avant et après \\
Couverture sur 88 incidents réels (red-team \#008) & de 75\,\% à 90\,\% des enregistrements, de 44\,\% à 81\,\% des
  commandes uniques \\
Temps du moteur sur les \nbexemples{} exemples (Linux, temps CPU, meilleur de trois) & médiane
  0,41\,ms ; 95\,\% sous 0,87\,ms ; maximum 1,70\,ms (budget : 10\,ms CPU) \\
Logique qui décide (champ \cmd{decided\_by} des reçus) & 116 lignes PowerShell du portail :
  0 décidée par les anciennes règles ; 806 lignes POSIX : 60, toutes des lignes qui nomment une
  cmdlet PowerShell ou un mot de \cmd{cmd.exe} (M5) \\
\bottomrule
\end{tabularx}
\end{table}

\paragraph{Les échecs connus.} Deux échecs connus restent hors du périmètre de l'outil : ils
dépendent de l'initialisation d'un dépôt git dans l'environnement de test. Les tests ignorés
sont ceux qui ne s'appliquent pas au système ou à l'environnement de la machine (tests propres à
Windows sous Linux, et inversement).

\paragraph{Ce que les chiffres ne disent pas.} Le nombre de tests mesure l'effort, pas la
sûreté. Les indicateurs qui comptent sont qualitatifs et vérifiables : chaque vulnérabilité
corrigée a un test qui échoue sans sa correction ; aucune partie n'a rendu une décision plus
permissive sans arbitrage ; et chaque décision affiche ce qu'elle n'a pas couvert.

